% SGA 6, Exposición 0: fragmento español alineado con la fuente.
% Alcance: secciones 5--7 y bibliografía completa, antes del Apéndice.
% Fuente francesa: sga6_fr_r1/src/fr.tex, líneas 417--546 inclusive.
% SHA-256 de la fuente francesa completa:
% 0D2158B9FE3B466309A92165727A0B191422E05FBAA0DACD5806E8322C54BA1A.

\subsection*{5.}
En la forma presentada en el n\textsuperscript{o}~4, la fórmula de
Riemann--Roch (1.2) también tiene sentido para un morfismo propio localmente
de intersección completa de espacios analíticos complejos (o de espacios
analíticos rígidos en el sentido de Tate). Pero, naturalmente, la
demostración dada en [RRR] solo se aplica a este caso cuando se supone,
además, que el morfismo $f$ es \underline{proyectivo} y, como en el caso
algebraico, parece que hace falta una idea esencialmente nueva para tratar el
caso general. La fórmula aquí considerada es una fórmula con valores en un
«~anillo de Chow analítico~» $\Gr^\bullet(X)$, o más bien
$\Gr^\bullet(X)_{\mathbb Q}$, construido según el principio explicado en
4.2 (la dificultad allí encontrada se desvanece, por otra parte, debido a
que estamos en característica cero, como se señala en Exp. XIV
n\textsuperscript{o}~1).

Cuando $X$ e $Y$ son no singulares, es posible dar otra versión de la fórmula
de Riemann--Roch procediendo como sigue. Sea primero $X$ un espacio analítico
complejo cualquiera, y sea $X^{\mathrm{top}}$ el mismo espacio provisto del
haz de funciones complejas continuas. Se tiene, pues, un morfismo canónico de
espacios anillados
\[
X^{\mathrm{top}}\longrightarrow X,
\]
que define, gracias al evidente carácter contravariante del funtor
$K^\bullet(X)$ respecto del espacio anillado variable $X$, un homomorfismo
de anillos
\begin{equation}\tag{5.1}
K^\bullet(X)\longrightarrow K^\bullet(X^{\mathrm{top}}).
\end{equation}
Supongamos, para simplificar, que $X^{\mathrm{top}}$ es compacto; entonces,
como en el caso (ya señalado en 3.1) de un esquema que admite un módulo
invertible amplio, se ve (Exp. II) que todo complejo perfecto sobre
$X^{\mathrm{top}}$ es globalmente isomorfo a un complejo acotado localmente
libre en cada grado, es decir, a un complejo de fibrados vectoriales
complejos sobre $X^{\mathrm{top}}$, y que, por consiguiente,
$K^\bullet(X^{\mathrm{top}})$ es canónicamente isomorfo al conocido grupo
$K(X^{\mathrm{top}})$ de los topólogos [1], definido en términos de fibrados
vectoriales complejos sobre el espacio compacto $X^{\mathrm{top}}$. Se puede
interpretar, pues, el homomorfismo (5.1) como un homomorfismo del
$K^\bullet(X)$ definido en términos de la estructura analítica al conocido
anillo $K(X^{\mathrm{top}})$ de los topólogos. Si ahora
$f:X\longrightarrow Y$ es un morfismo de espacios analíticos no singulares
compactos, se le sabe asociar, gracias a Atiyah--Hirzebruch [1], un
homomorfismo de grupos
\begin{equation}\tag{5.2}
f^{\mathrm{top}}_*:K^\bullet(X^{\mathrm{top}})\longrightarrow K^\bullet(Y^{\mathrm{top}}).
\end{equation}
Utilizando asimismo el homomorfismo $f_*$ de (3.1), definido por la fórmula
(3.2) (que aquí utiliza implícitamente el teorema de finitud de Grauert [\sisi{4}{3}] % Canon SGA-SRC-001493
para un morfismo propio de espacios analíticos), se obtiene un cuadrado de
homomorfismos naturales que solo hace intervenir grupos del tipo $K^\bullet$
(con exclusión de anillos del tipo de los anillos de Chow o de cohomología):
\begin{equation}\tag{5.3}
\begin{array}{ccc}
K^\bullet(X) & \longrightarrow & K^\bullet(X^{\mathrm{top}}) \\
f_*\downarrow && \downarrow f^{\mathrm{top}}_* \\
K^\bullet(Y) & \longrightarrow & K^\bullet(Y^{\mathrm{top}}).
\end{array}
\end{equation}
La fórmula de tipo Riemann--Roch a la que aludíamos consiste en afirmar la
conmutatividad de este cuadrado. Obsérvese la diferencia de naturaleza entre
esta afirmación (que no desprecia los fenómenos de torsión y se sitúa en un
grupo $K^\bullet(Y^{\mathrm{top}})$ de naturaleza esencialmente discreta) y
la fórmula de Riemann--Roch del tipo considerado en el
n\textsuperscript{o}~4, que se sitúa en un grupo del tipo «~Chow~» que ya no
es de naturaleza discreta, pero que, en cambio, debe tensorizarse por
$\mathbb Q$. Obsérvese que la fórmula (5.3) está demostrada, en todo caso,
cuando $X$ e $Y$ son variedades proyectivas no singulares [7], al igual que
la fórmula del tipo considerado en el n\textsuperscript{o}~4, que está
demostrada en este caso en [RRR]. Se trata de una variante plausible del
teorema del índice de Atiyah--Singer (y de sus generalizaciones debidas a
Shih [\sisi{7}{6}] y a Atiyah), que reclama evidentemente una generalización común. Una % Canon SGA-SRC-001494
vez demostrada la conmutatividad de (5.3), se deduciría inmediatamente,
mediante el «~teorema de Riemann--Roch diferenciable~» de
Atiyah--Hirzebruch [1], una variante cohomológica de la fórmula de
Riemann--Roch analítica compleja, a saber, la conmutatividad del cuadrado
\begin{equation}\tag{5.4}
\begin{array}{ccc}
K^\bullet(X) & \xrightarrow{\;x\longrightarrow \ch(x)\Todd(T_X)\;} & H^{2*}(X,\mathbb Q) \\
\downarrow && \downarrow \\
K^\bullet(Y) & \xrightarrow{\;y\longrightarrow \ch(y)\Todd(T_Y)\;} & H^{2*}(Y,\mathbb Q),
\end{array}
\end{equation}
donde las flechas horizontales se deducen de las clases de Chern
cohomológicas, deducidas del homomorfismo (5.1) y del homomorfismo de clase de
Chern ordinario
\[
K^\bullet(X^{\mathrm{top}})\longrightarrow H^{2*}(X,\mathbb Z),
\]
la primera flecha vertical es la misma que en (5.3), y la segunda es el
homomorfismo de Gysin en cohomología racional.

\subsection*{6.}
En todo lo anterior apenas se ha tratado más que el invariante contravariante
$K^\bullet(X)$, con exclusión del invariante covariante $K_\bullet(X)$ cuya
existencia se señaló en 3.1. Si se considera que $K^\bullet(X)$ proporciona
un análogo del anillo de \underline{cohomología} entera de un espacio
topológico, debe considerarse que $K_\bullet(X)$ proporciona un análogo de la
\underline{homología} entera, correspondiendo entonces la estructura de
módulo de $K_\bullet(X)$ sobre $K^\bullet(X)$ a la operación de producto
cap\footnote{Esta analogía no es, por otra parte, puramente verbal, y puede
precisarse, tras pasar a los graduados asociados $\Gr^\bullet(X)$,
$\Gr_\bullet(X)$, mediante una relación explícita de compatibilidad entre la
estructura de módulo sobre $\Gr_\bullet(X)$ y la operación «~cap~».}. El
hecho de que, cuando $X$ es regular, el homomorfismo natural
$K^\bullet(X)\to K_\bullet(X)$ sea un isomorfismo debe considerarse el
análogo del teorema de dualidad de Poincaré sobre las variedades topológicas
orientadas, que afirma que el producto cap con la clase de homología
fundamental define un isomorfismo de la cohomología a la homología. Estas
analogías sugieren, además, la conveniencia de definir también, para un
poliedro finito $X$, pongamos por caso, un invariante $K_\bullet(X)$ de
carácter covariante, cuyas relaciones con la homología entera (sucesión
espectral, homomorfismo de Chern\dots) serían análogas a las existentes entre
$K^\bullet(X)$ y la cohomología entera; tampoco parece excluido que puedan
existir, en términos de estos grupos $K_\bullet(X)$, variantes del teorema de
Riemann--Roch diferenciable, válidas para espacios más generales que las
variedades diferenciables compactas.

Volviendo al caso en que $X$ es un esquema noetheriano cualquiera, conviene
dotar a $K_\bullet(X)$ de una filtración creciente,
$\Fil_iK_\bullet(X)$ engendrada por clases $\cl_\bullet(F)$, con $F$ un haz
coherente sobre $X$ tal que $\dim \operatorname{supp}F\le i$. Se demostrará
(Exp. X) la compatibilidad de esta filtración con la estructura de módulo y
la filtración de $K^\bullet(X)$, es decir, la fórmula
\begin{equation}\tag{6.1}
\Fil^iK^\bullet(X)\cdot \Fil_jK_\bullet(X)\subset \Fil_{j-i}K_\bullet(X),
\end{equation}
lo que permite dotar al graduado asociado al módulo filtrado $K_\bullet(X)$,
denotado por
\[
\Gr_\bullet(X),
\]
de una estructura de módulo sobre $\Gr^\bullet(X)$. Sus elementos pueden
interpretarse, además, como \uline{clases de ciclos algebraicos} sobre $X$
(para una relación de equivalencia menos fina que la equivalencia racional,
cuando $X$ es de tipo finito sobre un cuerpo base $K$). Si
\[
f:X\longrightarrow Y
\]
es un morfismo propio de esquemas noetherianos, entonces
\begin{equation}\tag{6.2}
f_*:K_\bullet(X)\longrightarrow K_\bullet(Y)
\end{equation}
es compatible con las filtraciones y define un homomorfismo de grupos
graduados
\begin{equation}\tag{6.3}
f_*:\Gr_\bullet(X)\longrightarrow \Gr_\bullet(Y),
\end{equation}
que da lugar a una «~fórmula de proyección~», es decir, que es
$\Gr^\bullet(Y)$-lineal, fórmula que se obtiene al pasar a los graduados a
partir de la fórmula de proyección análoga para el homomorfismo (6.2) (que es
$K^\bullet(Y)$-lineal).

\uline{Desde el punto de vista propuesto en el Seminario, una teoría de las
intersecciones manejable sobre los esquemas noetherianos $X$, no
necesariamente regulares, debe consistir en el estudio combinado de los
invariantes contravariantes} $K^\bullet(X),\Gr^\bullet(X)$ \uline{y de los
invariantes covariantes} $K_\bullet(X),\Gr_\bullet(X)$, \uline{que desempeñan
papeles de naturaleza esencialmente diferente y cuyas propiedades se
complementan mutuamente.} Así, los grupos covariantes
$K_\bullet(X),\Gr_\bullet(X)$ se calculan mejor para ciertas fibraciones no
propias, gracias a la «~sucesión exacta de homotopía~» que puede establecerse
para ellos ([RRR] y Exp. IX); por ejemplo, se demuestran (Exp. IX)
isomorfismos canónicos tales como
\[
K_\bullet(X[t])\simeq K_\bullet(X),\qquad \Gr_\bullet(X[t])\simeq \Gr_\bullet(X),
\]
que fallan, en general, para $K^\bullet$ cuando $X$ es un esquema no regular.
En cambio, $K^\bullet(X)$ y $\Gr^\bullet(X)$ se prestan bien a los cálculos
gracias a sus estructuras de anillos.

Cuando se consideran esquemas $X$ propios sobre un cuerpo, la proyección
\[
f^X:X\longrightarrow(\operatorname{point})
\]
define un homomorfismo «~\uline{característica de Euler--Poincaré}~»
\begin{equation}\tag{6.4}
\lambda_X=f^X_*:K_\bullet(X)\longrightarrow K_\bullet(\operatorname{point})=\mathbb Z,
\end{equation}
que induce sobre el subgrupo $\Gr_0(X)=\Fil_0K_\bullet(X)$ el homomorfismo
«~grado de los $0$-ciclos~»
\begin{equation}\tag{6.5}
\deg_X:\Gr_0(X)\longrightarrow \mathbb Z.
\end{equation}
El homomorfismo de multiplicación
\[
\Gr^i(X)\times\Gr_i(X)\longrightarrow \Gr_0(X),
\]
compuesto con el homomorfismo de grado (6.5), proporciona entonces un
emparejamiento
\begin{equation}\tag{6.6}
\Gr^i(X)\times \Gr_i(X)\longrightarrow \mathbb Z,
\end{equation}
que explica que, en ciertos aspectos, $\Gr^\bullet(X)$ y $\Gr_\bullet(X)$
desempeñen papeles duales, al igual que la cohomología y la homología. Así,
si $f:X\to Y$ es un morfismo de esquemas algebraicos propios, los dos
homomorfismos
\[
f^*:\Gr^\bullet(Y)\longrightarrow \Gr^\bullet(X)\qquad\text{y}\qquad f_*:\Gr_\bullet(X)\longrightarrow \Gr_\bullet(Y)
\]
son formalmente transpuestos el uno del otro, en el sentido de los
emparejamientos (6.6).

Los homomorfismos (6.4) y (6.5), y el emparejamiento (6.6) que se deduce de
ellos, pueden considerarse la fuente de las \uline{relaciones numéricas} en
la teoría de las intersecciones sobre esquemas algebraicos propios sobre un
cuerpo dado $k$. Este punto de vista se desarrollará con algún detalle en
Exp. X, esencialmente independiente de Exp. VI, VII, VIII (centradas en la
demostración del teorema de Riemann--Roch) y de Exp. IX. Allí se darán
también algunas nociones sobre el \uline{comportamiento de}
$K^\bullet(X),\Gr^\bullet(X)$ \uline{por especialización}.

\subsection*{7.}
Como en toda teoría de intersección, importa precisar el lugar que ocupa en
la teoría el \uline{grupo de Picard} (o grupo de las \uline{clases de
divisores}, en la terminología clásica). Bajo una ligera restricción,
satisfecha automáticamente cuando $X$ es, por ejemplo, cuasicompacto, se
establece en Exp. X un isomorfismo canónico
\begin{equation}\tag{7.1}
c_1:\Pic(X)\xrightarrow{\sim}\Gr^1(X),
\end{equation}
que muestra el papel exacto del grupo de Picard en la teoría que proponemos.
Este papel justifica un estudio más detallado de dicho grupo de Picard en
las exposiciones XII y XIII, que son lógicamente independientes del texto
del Seminario. Kleiman realiza en Exp. XIII un estudio de la
\uline{equivalencia numérica} y de ciertos \uline{teoremas de finitud para el
grupo de Picard} (anunciados sin demostración en [\sisi{6}{5}]), utilizando técnicas de % Canon SGA-SRC-001495
carácter puramente numérico (que no utilizan la teoría del funtor de Picard
ni su representabilidad), debidas a él mismo y a Mumford, claramente más
sencillas que las demostraciones iniciales de Grothendieck. El único
resultado refractario a los métodos de Kleiman (y que se ve obligado a
utilizar) es [\sisi{6}{5}, C-07 (i)], cuya demostración y ciertos refinamientos
(incluidos algunos teoremas de representabilidad para el funtor de Picard)
ocupan la exposición XII. Esta última pertenece, por otra parte, a un
espíritu y una técnica completamente ajenos al resto del Seminario y
vinculados a [5]; esta exposición es lógicamente independiente de todas las
demás.

Finalmente, en Exp. XI, de carácter claramente más general que XII y XIII y
más próximo al espíritu general del Seminario, se estudia, sobre un topos
localmente anillado cualquiera, un funtor notable denominado
«~\uline{funtor determinante}~»:
\begin{equation}\tag{7.2}
\det^*:\underline{\Parf_{\mathrm{is}}}(X)\longrightarrow \underline{\Inv}(X)
\end{equation}
desde la categoría de los complejos perfectos sobre $X$, cuyos únicos
morfismos son los isomorfismos en el sentido de la categoría derivada $D(X)$,
hacia la categoría de los módulos invertibles sobre $X$. Este funtor está
descrito, salvo detalles menores, por la exigencia de ser de «~carácter
local~» y de venir dado, para un complejo $L$ acotado y con componentes
localmente libres, por la fórmula
\begin{equation}\tag{7.3}
\det^*(L_\bullet)\simeq\bigotimes_i \det(L_i)^{(-1)^i},
\end{equation}
donde $\det(L_i)$ designa la potencia exterior de grado máximo del módulo
localmente libre $L_i$. Aunque esta exposición también sea lógicamente
independiente del resto del Seminario, con excepción de la exposición I,
constituye, no obstante, un elemento importante del «~yoga~» general
propuesto en el Seminario y desarrollado esencialmente en las exposiciones I
a XI.

\section*{Bibliografía}
\begin{enumerate}[label={[\arabic*]},leftmargin=2.5em]
\item M. Atiyah y F. Hirzebruch, Vector bundles and homogeneous spaces, Symposia in Pure Mathematics, vol. 3, Differential Geometry, p. 7--38, 1961.
\item A. Borel y J.-P. Serre, Le th\'eor\`eme de Riemann--Roch, Bull. Soc. math. France, t. 86, p. 97--136 (1958).
\item H. Grauert, \emph{Ein Theorem der analytischen Garbentheorie}, Pub. Math. n\textsuperscript{o}~5, p. 5--64 (1960).
\item A. Grothendieck, Classes de faisceaux et th\'eor\`eme de Riemann--Roch, notas mimeografiadas, Princeton 1957 (citado [RRR] y reproducido como Apéndice a la Exp. 0 del Seminario).
\item A. Grothendieck, Technique de descente et th\'eor\`emes d'existence en G\'eom\'etrie Alg\'ebrique, en Fondements de la G\'eom\'etrie Alg\'ebrique, extractos del Séminaire Bourbaki 1957--1962, Secrétariat Mathématique, París.
\item W. Shih, Seminario de Topología en el IHES, 1966/67.
\item M. Atiyah y F. Hirzebruch, \emph{The Riemann--Roch theorem for analytic embeddings}, Topology, vol. 1, p. 151--166 (1962).
\end{enumerate}



\clearpage

